% TOP_PROBLEM: {"id": 3130, "title": "Signing a graph to have small magnitude eigenvalues", "category_id": 3, "category": {"id": 3, "name": "graph_theory", "display_name": "Graph Theory", "description": "Problems involving graphs, networks, and their properties.", "slug": "graph-theory", "order_index": 3, "created_at": "2026-07-31T15:26:25.671Z"}, "status": "open", "problem_number": "OPG-48264", "set_id": 12, "statusQualification": "The necessary degree range d at least two is explicit; d=1 would give an immediate counterexample to the unqualified shorthand."}
% FIRST_POSTED: 2026-10-01
% ATTEMPT_STATUS: unresolved
% UnsolvedMath ID 3130; source catalogue code OPG-48264
% !TeX program = lualatex
% GPT-6 Astra Ultra research attempt, 2026-10-01; self-contained partial work.
% Completion estimate: no numerical percentage assigned; full problem unresolved.
\documentclass[11pt,a4paper]{article}
\usepackage[margin=25mm]{geometry}
\usepackage{fontspec}
\setmainfont{lmroman10-regular.otf}[BoldFont=lmroman10-bold.otf,
 ItalicFont=lmroman10-italic.otf,BoldItalicFont=lmroman10-bolditalic.otf]
\usepackage{amsmath,amssymb,mathtools}
\usepackage{unicode-math}
\setmathfont{latinmodern-math.otf}
\AtBeginDocument{\let\setminus\smallsetminus}
\usepackage{xurl,hyperref}
\hypersetup{colorlinks=true,urlcolor=blue}
\setcounter{secnumdepth}{0}
\setlength{\parindent}{0pt}
\setlength{\parskip}{5pt}
\setlength{\emergencystretch}{3em}
\begin{document}
\section*{Signing a graph to have small magnitude eigenvalues}\label{3130}
{\small UnsolvedMath ID 3130 · OPG-48264 · Graph Theory · OpenGarden}
\subsection{Definitions and mathematical statement}
Let $d\ge2$ and let $G=(V,E)$ be a finite simple undirected $d$-regular
graph, meaning that every vertex has exactly $d$ neighbours.
For a signing $s:E\to\{-1,1\}$ define the real symmetric matrix
$A_s$, indexed by $V$, by
\[
 (A_s)_{uv}=\begin{cases}s(\{u,v\}),&\{u,v\}\in E,\\
                         0,&\text{otherwise}.\end{cases}
\]
In particular its diagonal is zero, and the two entries corresponding
to the same undirected edge receive the same sign. A signing does not
change which entries of the adjacency matrix are zero.

The conjecture asserts that for every such $G$ there is a signing with
\[
 \max_{\lambda\in\operatorname{Spec}(A_s)}|\lambda|
                 \le2\sqrt{d-1}.
\]
All eigenvalues are real because $A_s$ is symmetric. The inequality
bounds both the largest eigenvalue and the magnitude of the most
negative one; a bound on only one end of the spectrum is insufficient.
Equivalently, the Euclidean operator norm of $A_s$ must have this bound.

Every eigenvalue of the signed matrix is included, with no exception
for eigenvectors related to those of the unsigned regular graph.
The signs may depend on the whole graph and need not come from
assigning signs to vertices. No bipartiteness or connectedness is imposed.
The source's implicit range is specified as $d\ge2$: at $d=1$ a single
edge has signed eigenvalues $1,-1$, whereas the proposed right-hand
side is zero, so that degenerate case cannot be included.
\subsection{Short English statement}
Can every regular graph of degree at least two be signed so that all signed adjacency eigenvalues have magnitude at most $2\sqrt{d-1}$?
\subsection{Research attempt}
\paragraph{Scope and process.}
GPT-6 Astra Ultra, 1 October 2026. This attempt keeps the catalogue statement above, checks primary literature, and develops a restricted argument with an adversarial gap audit. The proofs below are the mathematical output of this attempt, not a claim of novelty or external certification.

\paragraph{Literature and attack plan.}
The Bilu--Linial framework is cited in [S3]. Marcus, Spielman and Srivastava [R1] obtain a signing controlling the largest eigenvalue; bipartite symmetry converts that result to the required two-sided estimate on bipartite graphs. Their one-sided statement must not be silently promoted to the two-sided nonbipartite conjecture. Here we prove the symmetry implication, settle the degree-two case directly, construct a nonbipartite example, and test two tempting but invalid shortcuts.

\paragraph{The degree-two case needs no sign choice.}
For any graph of maximum degree at most $d$, any signing, and any real vector $x$,
\[
 |x^{\mathsf T}A_sx|
 \le2\sum_{uv\in E}|x_ux_v|
 \le\sum_{uv\in E}(x_u^2+x_v^2)
 \le d\sum_vx_v^2. \tag{1}
\]
As $A_s$ is symmetric, the Rayleigh-quotient bound gives $\|A_s\|_2\le d$. For $d=2$ this equals $2\sqrt{d-1}$, proving the target for every 2-regular graph and indeed every signing. For $d>2$, however,
\[
 d^2-4(d-1)=(d-2)^2>0,
\]
so the same estimate is strictly too weak.

\paragraph{What bipartiteness supplies.}
For a bipartition $V=X\cup Y$, let $J$ be diagonal, with entry $+1$ on $X$ and $-1$ on $Y$. Since every edge crosses the partition, entrywise multiplication gives
\[
 JA_sJ=-A_s.
\]
Thus if $A_sx=\lambda x$, then $A_s(Jx)=-\lambda(Jx)$. The spectrum is symmetric, and bounding its largest eigenvalue bounds its whole magnitude. Applying the one-sided signing theorem of [R1] therefore proves the desired existence statement for the bipartite regular subclass. The deep existence theorem is cited, not reproved; the passage from one-sided to two-sided control has been proved here with its exact hypothesis.

\paragraph{An explicit nonbipartite signing.}
For $K_4$, make the edge $12$ negative and all other edges positive. The signed adjacency matrix is
\[
 A_s=\begin{pmatrix}
 0&-1&1&1\\-1&0&1&1\\1&1&0&1\\1&1&1&0
 \end{pmatrix}.
\]
The vector $(1,-1,0,0)$ has eigenvalue $1$, and $(0,0,1,-1)$ has eigenvalue $-1$. On the remaining invariant subspace of vectors $(a,a,b,b)$, the action is represented by
\[
 \begin{pmatrix}-1&2\\2&1\end{pmatrix},
\]
whose square is $5I$. Its eigenvalues are $\sqrt5,-\sqrt5$. The full spectrum is therefore $\{1,-1,\sqrt5,-\sqrt5\}$ and its norm is $\sqrt5<2\sqrt2$, as required for degree three. This directly checks a genuinely nonbipartite example, without a numerical eigenvalue approximation.

\paragraph{Why vertex signs alone cannot work.}
If one assigns signs $t_v\in\{-1,1\}$ to vertices and sets $s(uv)=t_ut_v$, then $A_s=TAT$, where $T=\operatorname{diag}(t_v)$ and $T^{-1}=T$. Hence the signed and unsigned matrices have identical spectra. A $d$-regular graph has unsigned eigenvalue $d$ on the all-ones vector, so such a signing has norm at least $d$. By the strict inequality after (1), vertex-induced signs cannot meet the target for any $d>2$. Edge signs must carry information not removable by this diagonal conjugation.

\paragraph{Why flipping all signs does not combine the two bounds.}
Replacing $s$ by $-s$ negates the matrix and exchanges its extreme eigenvalues. A signing good at the upper end can therefore give another signing good at the lower end; this is not one signing good at both ends. On $K_4$, the all-negative signing has spectrum $\{-3,1,1,1\}$. Its largest eigenvalue is comfortably below $2\sqrt2$, while its norm is three and fails the target. This is an explicit obstruction to inferring the two-sided bound from an upper-end certificate alone.

\paragraph{Verification, gap, and final status.}
The displayed four eigenvectors/subspaces account for all four dimensions. The $K_4$ example respects a common sign for the two matrix entries of each undirected edge and has zero diagonal. For disconnected regular graphs, matrices are block diagonal, so each component can be signed separately. Degree one remains excluded exactly as in the catalogue.

The unresolved step is to choose a single edge signing of every nonbipartite regular graph of degree at least three that controls both spectral ends simultaneously. The degree-two proof, cited bipartite theorem, and explicit example do not establish that choice. The full problem is \textbf{UNRESOLVED} by this attempt.

\subsection{Sources}
\begin{itemize}
\item[S1] ULAM.AI and the UnsolvedMath Contributors, supplied problems.json, record 3130 (OPG-48264), OpenGarden collection. Catalogue identity and source transcription; CC BY 4.0.
\url{https://huggingface.co/datasets/ulamai/UnsolvedMath}.
\item[S2] Open Problem Garden, Signing a graph to have small magnitude eigenvalues. Original problem and discussion; the mathematical formulation below is expanded from the supplied source record.
\url{http://www.openproblemgarden.org/op/signing_a_graph_to_have_small_magnitude_eigenvalues}.
\item[S3] Y. Bilu and N. Linial, Constructing expander graphs by 2-lifts and discrepancy vs. spectral gap, original signing framework, arXiv:math/0312022.
\url{https://arxiv.org/abs/math/0312022}.
\item[R1] A. W. Marcus, D. A. Spielman and N. Srivastava, \emph{Interlacing Families I: Bipartite Ramanujan Graphs of All Degrees}, arXiv:1304.4132; Annals of Mathematics 182 (2015), 307--325. Primary paper checked for the one-sided signing theorem and bipartite consequence.
\url{https://arxiv.org/abs/1304.4132}.
\url{https://www.cs.yale.edu/homes/spielman/PAPERS/Marcus_Spielman_SrivastavaIFI.pdf}.
\end{itemize}
\end{document}
